\documentclass[10pt]{beamer}
\usetheme{metropolis}

\usepackage{amsmath,amssymb}
\usepackage{mathtools}
\usepackage{xeCJK}
\setCJKmainfont{Noto Serif SC}
\usepackage{tikz-cd}

% 1. 开启标题背景色填充
\metroset{block=fill}

% theorem environments
\newtheorem{proposition}{Proposition}
\newtheorem{exercise}{Exercise}

% notation
\newcommand{\R}{\mathbb{R}}
\newcommand{\N}{\mathbb{N}}
\newcommand{\F}{\mathcal{F}}
\newcommand{\supp}{\operatorname{supp}}

\title{Ch2\quad Tangent Vectors and Differentials}
\subtitle{Germs of Functions \& The Tangent Space}
\date{September 28, 2026}
\author{LXJ}

\begin{document}

\maketitle

%====================================================================
\section{The Idea: Derivations}

\begin{frame}{Tangent vectors as operators}
	\textbf{Idea.} Any vector $v = (v_1, \dots, v_n) \in \R^n$ can be thought
	of as an \emph{operator} on differentiable functions. If $f$ is
	differentiable at $p$, then $v$ assigns to $f$ the
	\emph{directional derivative}
	\[
		v(f) = \sum_{i=1}^{n} v_i \, \frac{\partial f}{\partial r_i}(p),
		\qquad \text{which is a real value.}
	\]
	It satisfies
	\begin{enumerate}
		\item \textbf{linearity}\quad
		      $v(\lambda f + \mu g) = \lambda\, v(f) + \mu\, v(g)$;
		\item \textbf{Leibniz rule}\quad
		      $v(f \cdot g) = f(p)\cdot v(g) + v(f)\cdot g(p)$.
	\end{enumerate}
\end{frame}

\begin{frame}{Why germs of functions?}
	The operation of taking derivatives only depends on the
	\emph{local} properties of functions. This is captured by the concept
	of \emph{germs of functions}.
	\begin{block}{Remark}
		A vector $v$ is thus ``defined'' as a linear function on
		differentiable functions, determined by the values of $v(f)$.
	\end{block}
	\begin{definition}[Germs of functions]
		Let $m \in M$. Functions $f$ and $g$ defined on open sets
		containing $m$ are said to have the \emph{same germ at $m$} if they
		agree on some neighborhood of $m$. The equivalence classes of germs
		of $C^\infty$ functions are called \emph{germs}, denoted by
		$\widetilde{\F}_m$, i.e.
		\[
			\widetilde{\F}_m = \bigl\{ [f] \bigm|
			f, g \in [f] \ \text{if } \exists \text{ open set } U \ni m,
			\ \text{s.t.\ } f\big|_U = g\big|_U \bigr\}.
		\]
	\end{definition}
\end{frame}

%====================================================================
\section{The Algebra of Germs}

\begin{frame}{Operations on $\widetilde{\F}_m$}
	The addition, scalar multiplication and multiplication of functions
	induce
	\begin{align*}
		\text{(addition)}       & \quad [f] + [g] = [f + g],       \\
		\text{(scalar mult.)}   & \quad \lambda [f] = [\lambda f], \\
		\text{(multiplication)} & \quad [f] \cdot [g] = [fg].
	\end{align*}
	\begin{exercise}
		Prove that these are well-defined operations.
	\end{exercise}
	\begin{block}{Remark}
		These operations make $\widetilde{\F}_m$ an \emph{algebra} over
		$\R$. A germ $[f]$ has a well-defined value $f(m)$ at $m$.
	\end{block}
	\begin{definition}[The ideals $F_m$ and $F_m^k$]
		\[
			F_m = \bigl\{ [f] \bigm| f \in \widetilde{\F}_m,\ f(m) = 0 \bigr\}.
		\]
	\end{definition}
\end{frame}

\begin{frame}{Powers of the ideal $F_m$}
	$F_m$ is an \emph{ideal} of $\widetilde{\F}_m$. Let $F_m^k$ be the
	$k$-th power of $F_m$: the ideal of $\widetilde{\F}_m$ consisting of
	all finite linear combinations of $k$-fold products of elements of
	$F_m$,
	\[
		F_m^k = \bigl\{ \text{linear comb.\ of } [f_1] \cdots [f_k],
		\ [f_i] \in F_m \bigr\}.
	\]
	These form a descending sequence of ideals
	\[
		\widetilde{\F}_m \supset F_m \supset F_m^2 \supset F_m^3
		\supset \cdots
	\]
\end{frame}

%====================================================================
\section{The Tangent Space}

\begin{frame}{Tangent vectors as derivations}
	\begin{definition}[Tangent vector]
		A \emph{tangent vector} $v$ at the point $m \in M$ is a
		\emph{linear derivation} of the algebra $\widetilde{\F}_m$, i.e.\
		for $[f], [g] \in \widetilde{\F}_m$ and $\lambda, \mu \in \R$:
		\begin{enumerate}
			\item[(a)] $v\bigl(\lambda [f] + \mu [g]\bigr)
				      = \lambda\, v([f]) + \mu\, v([g])$;
			\item[(b)] $v\bigl([f]\cdot[g]\bigr)
				      = v([f])\, g(m) + f(m)\, v([g])$.
		\end{enumerate}
	\end{definition}
	Let $T_m M$ or $M_m$ denote the set of all linear derivations at $m$,
	called the \emph{tangent space to $M$ at $m$}.
	\begin{block}{Note}
		If $M_m \neq \varnothing$, then $M_m$ is a \emph{linear space},
		since for $v, w \in M_m$, $\lambda, \mu \in \R$,
		\[
			(\lambda v + \mu w)([f])
			\ \triangleq\ \lambda\, v([f]) + \mu\, w([f]).
		\]
	\end{block}
\end{frame}

\begin{frame}{Two remarks on the definition}
	\begin{enumerate}
		\item This definition of tangent spaces only works for the
		      $C^\infty$ case;
		\item It is \emph{intrinsic} --- it does not depend on the local
		      coordinate system.
	\end{enumerate}
	\medskip
	\begin{block}{Note (constants are killed)}
		For a derivation $w \in M_m$ and $[c]$ the germ of a locally
		constant function, $w([c]) = 0$.
	\end{block}
	\begin{proof}
		It suffices to show $w([1]) = 0$:
		\[
			w([1]) = w([1]\cdot[1])
			= w([1]) \cdot 1 + 1 \cdot w([1]) = 2\, w([1]). \qedhere
		\]
	\end{proof}
\end{frame}

%====================================================================
\section{$M_m \cong (F_m / F_m^2)^*$}

\begin{frame}{The key isomorphism}
	\begin{lemma}
		$M_m$ is isomorphic to $(F_m / F_m^2)^*$.
	\end{lemma}
	\begin{proof}
		Consider $v \in M_m$. Then $v$ is a linear function on $F_m$ by
		$v : [f] \mapsto v([f])$. This function \emph{vanishes on $F_m^2$}
		by the Leibniz rule of $v$: $\forall\, [f], [g] \in F_m$,
		\[
			v([f][g]) = v([f]) \cdot g(m) + f(m) \cdot v([g]) = 0.
		\]
		So $v$ can be thought of as a linear function defined on the
		cosets $[f] + F_m^2$, therefore $v \in (F_m/F_m^2)^*$.
	\end{proof}
\end{frame}

\begin{frame}{Converse direction}
	Conversely, for $\ell \in (F_m / F_m^2)^*$, we can define
	$v_\ell \in M_m$ by
	\[
		v_\ell([f]) \ \triangleq\
		\ell\bigl( [f] - [f(m)] + F_m^2 \bigr).
	\]
	Then one verifies that $v_\ell$ is a derivation on
	$\widetilde{\F}_m$:
	\begin{enumerate}
		\item \textbf{Linearity.}
		      $v_\ell\bigl(\lambda[f] + \mu[g]\bigr)
			      = \ell\bigl( \lambda[f] - \lambda[f(m)] + F_m^2 \bigr)
			      + \ell\bigl( \mu[g] - \mu[g(m)] + F_m^2 \bigr)
			      = \lambda\, v_\ell([f]) + \mu\, v_\ell([g])$.
		\item \textbf{Leibniz.}
		      $v_\ell\bigl([f][g]\bigr)
			      = \ell\bigl( [f][g] - [f(m)][g(m)] + F_m^2 \bigr)$,
		      where
		      \begin{align*}
			      [f][g] - [f(m)][g(m)]
			       & = \underbrace{\bigl([f]-[f(m)]\bigr)\bigl([g]-[g(m)]\bigr)}_{\in\, F_m^2} \\
			       & \quad + [f(m)]\bigl([g]-[g(m)]\bigr)
			      + \bigl([f]-[f(m)]\bigr)[g(m)],
		      \end{align*}
		      so $v_\ell([f][g]) = f(m)\, v_\ell([g]) + v_\ell([f])\, g(m)$.
	\end{enumerate}
\end{frame}

\begin{frame}{An exercise hidden in the proof}
	\begin{exercise}
		Prove that for $\ell \in (F_m / F_m^2)^*$,
		\[
			\ell\bigl( [f(m)]\cdot\bigl([g]-[g(m)]\bigr) \bigr)
			= f(m)\, \ell\bigl( [g] - [g(m)] \bigr).
		\]
	\end{exercise}
	\bigskip
	\begin{alertblock}{Theorem}
		\[
			\dim\bigl(F_m / F_m^2\bigr) = \dim M.
		\]
	\end{alertblock}
\end{frame}

%====================================================================
\section{Proof of $\dim(F_m/F_m^2) = \dim M$}

\begin{frame}{A Taylor lemma with integral remainder}
	\small
	\begin{lemma}
		If $g \in C^k(U)$ for $k \geq 2$ and $U$ is a convex open set in
		$\R^d$, then $\forall\, p, q \in U$,
		\begin{multline*}
			g(q) = g(p) + \sum_{i=1}^{d}
			\left. \frac{\partial g}{\partial r_i} \right|_{p}
			\bigl(r_i(q) - r_i(p)\bigr) \\
			+ \sum_{i,j} \bigl(r_i(q) - r_i(p)\bigr)
			\bigl(r_j(q) - r_j(p)\bigr)
			\int_0^1 (1-t)
			\left. \frac{\partial^2 g}{\partial r_j \partial r_i}
			\right|_{p + t(q - p)} dt .
		\end{multline*}
	\end{lemma}
	\begin{proof}
		Since $U$ is convex, for any $p, q \in U$ the segment
		$\{ p + t(q-p) \mid t \in [0,1] \} \subset U$. Writing
		$p_i \triangleq r_i(p)$, $q_i \triangleq r_i(q)$, we get
		\[
			g(q) - g(p)
			= \int_0^1 \frac{d}{dt}\, g\bigl(p + t(q-p)\bigr)\, dt
			= \sum_i \int_0^1
			\left. \frac{\partial g}{\partial r_i}
			\right|_{p + t(q-p)} (q_i - p_i)\, dt .
			\qedhere
		\]
	\end{proof}
\end{frame}

\begin{frame}{Proof of the Taylor lemma (continued)}
	For each integrand, expand around $p$:
	\[
		\left. \frac{\partial g}{\partial r_i}
		\right|_{p + t(q-p)}
		= \left. \frac{\partial g}{\partial r_i} \right|_{p}
		+ \sum_j \int_0^t
		\frac{\partial^2 g}{\partial r_j \partial r_i}
		\bigl(p + s(q-p)\bigr)\, (q_j - p_j)\, ds .
	\]
	So, interchanging the order of integration ($s \leftrightarrow t$),
	\begin{align*}
		g(q) & = g(p) + \sum_i
		\left. \frac{\partial g}{\partial r_i} \right|_{p} (q_i - p_i)  \\
		     & \quad + \sum_{i,j} (q_i - p_i)(q_j - p_j) \int_0^1 (1-t)
		\left. \frac{\partial^2 g}{\partial r_j \partial r_i}
		\right|_{p + t(q-p)} dt .
		\qedhere
	\end{align*}
\end{frame}

\begin{frame}{Proof of the theorem: spanning set}
	Let $(U, \varphi)$ be any coordinate system about $m \in M$, with
	coordinate functions $(x_1, \dots, x_d)$, $\mathbf{x}(p) =(u_1,\cdots, u_d)$, $d = \dim M$. Let
	$[f] \in \widetilde{\F}_m$. Applying $f \circ \varphi^{-1}$ to the
	lemma:
	\begin{multline*}
		f \circ \varphi^{-1}(\mathbf{u}) - f \circ \varphi^{-1}(\mathbf{u}_0)
		= \sum_{i=1}^{d}
		\left. \frac{\partial (f \circ \varphi^{-1})}{\partial u_i}
		\right|_{\mathbf{u}_0} (u_i - u_i^0) \\
		+ \sum_{i,j} (u_i - u_i^0)(u_j - u_j^0) \int_0^1 (1-t)
		\left. \frac{\partial^2 (f \circ \varphi^{-1})}
		{\partial u_j \partial u_i} \right|_{\mathbf{u}_0 + t(\mathbf{u} - \mathbf{u}_0)} dt .
	\end{multline*}
	Since $\mathbf{x}(p) := \varphi(p)=\mathbf{u}$, $x_0 = \varphi(m)$, composing with $\varphi$
	gives, \emph{modulo} $F_m^2$,
	\[
		[f(p)] - [f(m)]
		= \sum_{i=1}^{d}
		\left. \frac{\partial (f \circ \varphi^{-1})}{\partial u_i}
		\right|_{\varphi(m)} \bigl(x_i(p) - x_i(m)\bigr).
	\]
	Hence $\bigl\{ [x_i] - [x_i(m)] \bigr\}$ \emph{spans} $F_m / F_m^2$.
\end{frame}

\begin{frame}{Proof of the theorem: linear independence}
	We claim they are also \emph{linearly independent}. Suppose that
	\[
		\sum_i a_i \bigl(x_i - x_i(m)\bigr) \in F_m^2
		\ \Longleftrightarrow\
		\sum_i a_i \bigl(u_i - u_i^0\bigr) \in F_{\varphi(m)}^2 .
	\]
	Then applying $\left. \dfrac{\partial}{\partial u_j}
		\right|_{\varphi(m)}$ to both sides gives
	\[
		a_j = 0 \qquad \forall\, j .
	\]
	Therefore $\bigl\{ [x_i] - [x_i(m)] \bigr\}_{i=1}^{d}$ is a basis of
	$F_m / F_m^2$, and
	\[
		\dim\bigl(F_m / F_m^2\bigr) = d = \dim M . \qed
	\]
\end{frame}

%====================================================================
\section{The Usual View: Coordinate Basis}

\begin{frame}{Tangent vectors acting on functions}
	In practice, a tangent vector operates on \emph{functions} rather than
	on their germs. Let $f$ be a differentiable function defined on a
	neighborhood of $p \in M$. For $v \in T_p M$ (or $M_p$), we define
	\[
		v(f) \ \triangleq\ v([f]).
	\]
	So $v(f) = v(g)$ whenever $f$ and $g$ agree on a neighborhood of $p$.
	\[
		v(\lambda f + \mu g) = \lambda\, v(f) + \mu\, v(g)
		\tag{a}
	\]
	\[
		v(f \cdot g) = v(f)\, g(p) + f(p)\, v(g)
		\tag{b}
	\]
\end{frame}

\begin{frame}{Coordinate tangent vectors}
	\begin{definition}[Coordinate tangent vectors]
		Let $(U, \varphi)$ be a coordinate system with coordinate functions
		$x_1, \dots, x_d$, and let $p \in U$. For $i = 1, \dots, d$ define
		\[
			\left. \frac{\partial}{\partial x_i} \right|_{p} (f)
			\ \triangleq\
			\left. \frac{\partial (f \circ \varphi^{-1})}{\partial r_i}
			\right|_{\varphi(p)}
		\]
		for each $C^\infty$ function $f$ on a neighborhood of $p$. We also
		use the notation
		$\left. \frac{\partial f}{\partial x_i} \right|_{p}
			= \left. \frac{\partial}{\partial x_i} \right|_{p} (f)$.
	\end{definition}
	\begin{block}{Remark (a)}
		Clearly $\left. \frac{\partial}{\partial x_i} \right|_{p}(f)$
		depends only on the germ of $f$ at $p$, and (a), (b) are satisfied.
		So $\left. \frac{\partial}{\partial x_i} \right|_{p}$ is a tangent
		vector at $p$.
	\end{block}
\end{frame}

\begin{frame}{The coordinate basis}
	Moreover
	$\Bigl\{ \left. \frac{\partial}{\partial x_i} \right|_{p}
		\Bigr\}_{i=1}^{d}$ is a basis of $T_p M$ --- indeed the basis
	\emph{dual} to $\bigl\{ [x_i - x_i(p)] \bigr\}$ of $F_p / F_p^2$,
	since
	\[
		\left. \frac{\partial}{\partial x_i} \right|_{p}
		\bigl(x_j - x_j(p)\bigr) = \delta_{ij}.
	\]
	\begin{enumerate}
		\item[(b)] If $w \in T_p M$, then
		      \[
			      w = \sum_{i=1}^{d} w(x_i)
			      \left. \frac{\partial}{\partial x_i} \right|_{p},
		      \]
		      since both sides give the same result when acting on the
		      functions $x_i - x_i(p)$.
	\end{enumerate}
\end{frame}

\begin{frame}{Change of coordinates}
	\begin{enumerate}
		\item[(c)] If $(U, \varphi)$ and $(V, \psi)$ are coordinate systems
		      about $p$ with coordinate functions $(x_1, \dots, x_d)$ and
		      $(y_1, \dots, y_d)$, then by (b)
		      \[
			      \left. \frac{\partial}{\partial y_j} \right|_{p}
			      = \sum_{i=1}^{d} \left. \frac{\partial x_i}{\partial y_j}
			      \right|_{p} \left. \frac{\partial}{\partial x_i} \right|_{p}.
		      \]
		      Note $\bigl( \frac{\partial}{\partial x_i} \bigr)$ depends on
		      $\varphi$, not only on $x_i$; so $x_i = y_i$ does
		      \emph{not} necessarily imply
		      $\frac{\partial}{\partial x_i} = \frac{\partial}{\partial y_i}$.
		\item[(d)] On the Euclidean space $\R^d$, the tangent vectors are
		      all the linear combinations of $\frac{\partial}{\partial r_i}$.
	\end{enumerate}
\end{frame}

\begin{frame}{The $C^k$ caveat}
	\begin{block}{Note}
		The proof that $\dim (F_m / F_m^2) = \dim M$ fails in the $C^k$
		case ($k < \infty$): the remainder term in the Taylor lemma is only
		$C^{k-1}$, and $F_m / F_m^2$ is in fact always infinite-dimensional
		in the $C^k$ case.
	\end{block}
	There are various ways to define tangent vectors in the $C^k$ case so
	that $\dim T_p M = \dim M$. For example, one may define a tangent
	vector at $p$ as a mapping assigning to each function ($C^k$ locally
	about $p$) a real number such that, whenever $(U, \varphi)$ is a
	coordinate system on a neighborhood of $p$, there exists a list of
	numbers $(a_1, \dots, a_d)$ with
	\[
		v(f) = \sum_{i=1}^{d} a_i
		\left. \frac{\partial (f \circ \varphi^{-1})}{\partial r_i}
		\right|_{\varphi(p)} .
	\]
\end{frame}

%====================================================================
\section{The Differential}

\begin{frame}{The differential of a map}
	\begin{definition}[The differential]
		Let $\Psi : M \to N$ be a $C^\infty$ map and $p \in M$. The
		\emph{differential} (or tangent map) of $\Psi$ at $p$ is the linear
		map
		\[
			\Psi_{*, p} : T_p M \longrightarrow T_{\Psi(p)} N
		\]
		defined as follows: for $w \in T_p M$ and a $C^\infty$ function
		$f$ on a neighborhood of $\Psi(p)$,
		\[
			\Psi_{*, p}(w)(f) \ \triangleq\ w(f \circ \Psi)
			= d\Psi_p(w)(f).
		\]
	\end{definition}
	$\Psi$ is called \emph{non-singular} at $p$ if
	$\ker (\Psi_{*, p}) = \{0\}$.
	\medskip

	The dual map
	$\Psi_p^* : (T_{\Psi(p)} N)^* \longrightarrow (T_p M)^*$ is defined by
	$\Psi_p^*(\omega)(w) = \omega(\Psi_*(w)) = \omega(d\Psi(w))$.
\end{frame}

\begin{frame}{Example: the differential of a function}
	If $f : M \to \R$ is smooth, $w \in T_p M$ and $f(p) = r_0$, then
	\[
		df(w) = w(f) \left. \frac{d}{dr} \right|_{r_0}.
	\]
	\begin{proof}
		If $w = \sum_i v_i \left. \frac{\partial}{\partial x_i} \right|_p$,
		then $df(w) = \sum_i v_i\, df\bigl( \frac{\partial}{\partial x_i} \bigr)$.
		For any $C^\infty$ function $g$ about $r_0$,
		\[
			df(w)(g) = w(g \circ f)
			= \sum_i v_i \frac{\partial g}{\partial r}
			\frac{\partial f}{\partial x_i}
			= w(f)\, \frac{d}{dr}(g).
			\qedhere
		\]
	\end{proof}
	In this case $df \in (T_p M)^*$ is simply given by $df(w) \triangleq w(f)$.
\end{frame}

\begin{frame}{Jacobian matrix and dual basis}
	\begin{enumerate}
		\item[(a)] Let $(U, x_1, \dots, x_d)$ and $(V, y_1, \dots, y_d)$ be
		      coordinate systems about $p$ and $\Psi(p)$. Then
		      \[
			      d\Psi\left( \left. \frac{\partial}{\partial x_j} \right|_{p} \right)
			      = \sum_{i=1}^{d} \left. \frac{\partial (y_i \circ \Psi)}
			      {\partial x_j} \right|_{p}
			      \left. \frac{\partial}{\partial y_i} \right|_{\Psi(p)} .
		      \]
		      The matrix $\bigl( \left. \frac{\partial (y_i \circ \Psi)}
			      {\partial x_j} \right|_{p} \bigr)$ is called the
		      \emph{Jacobian} of $\Psi$.
		\item[(b)] If $(U, x_1, \dots, x_d)$ is a coordinate system on $M$
		      and $p \in U$, then $\bigl\{ (dx_i)_p \bigr\}$ is the basis of
		      $(T_p M)^*$ dual to
		      $\bigl\{ \left. \frac{\partial}{\partial x_i} \right|_{p} \bigr\}$.
		      For $f : M \to \R$ smooth,
		      \[
			      df_p = \sum_{i=1}^{d} \left. \frac{\partial f}{\partial x_i}
			      \right|_{p} (dx_i)_p .
		      \]
	\end{enumerate}
\end{frame}

\begin{frame}{Chain rule and pullback}
	\begin{enumerate}
		\item[(c)] \textbf{Chain rule.} Let $\Psi : M \to N$ and
		      $\Phi : N \to X$ be $C^\infty$ maps. Then
		      \[
			      d(\Phi \circ \Psi)_p
			      = d\Phi\big|_{\Psi(p)} \circ\, d\Psi\big|_{p} .
		      \]
		\item[(d)] If $\Psi : M \to N$ and $f : N \to \R$ are $C^\infty$,
		      then
		      \[
			      \Psi^*(df)_{\Psi(p)} = d(f \circ \Psi)_p,
		      \]
		      because for all $w \in T_p M$,
		      \[
			      \Psi^*(df)_{\Psi(p)}(w) = d(f \circ \Psi)\big|_{p}(w).
		      \]
	\end{enumerate}
\end{frame}

%====================================================================
\section{Smooth Curves}

\begin{frame}{Smooth curves}
	\begin{definition}[Smooth curve]
		A $C^\infty$ map $\sigma : (a, b) \to M$ is called a \emph{smooth
			curve} in $M$. Its tangent vector at $t \in (a,b)$ is
		\[
			\dot{\sigma}(t) \ \triangleq\
			d\sigma\left( \left. \frac{d}{dt} \right|_{t} \right)
			= \sigma_*\left( \left. \frac{d}{dt} \right|_{t} \right)
			\in M_{\sigma(t)},
		\]
		given in coordinates by
		$\dot{\sigma}(t) = \sum_{i=1}^{d} \left. \frac{\partial x_i}{\partial t} \right|_{t} \left. \frac{\partial}{\partial x_i} \right|_{\sigma(t)}$.
	\end{definition}
	\begin{block}{Remark}
		Every $w \in T_p M$ is the tangent vector to a smooth curve: choose
		$(U, \varphi)$ centered at $p$ with
		$w = d\varphi^{-1}\bigl( \left. \frac{\partial}{\partial r_1}
			\right|_{0} \bigr)$; then $w$ is tangent at $0$ to
		$t \mapsto \varphi^{-1}(t, 0, \dots, 0)$.
	\end{block}
\end{frame}

\begin{frame}{Remarks on curves}
	\begin{itemize}
		\item Two curves $\sigma, \tau$ with $\sigma(t_0) = \tau(t_0) = p$
		      have the same tangent vector at $t_0$ iff
		      $\left. \frac{d(f \circ \sigma)}{dt} \right|_{t_0}
			      = \left. \frac{d(f \circ \tau)}{dt} \right|_{t_0}$ for all
		      $f$ which are $C^\infty$ about $p$.
		\item In $\R^d$:\quad
		      $\dot{\sigma}(t) = \sum_i \left. \frac{d\sigma_i}{dt}
			      \right|_{t} \left. \frac{\partial}{\partial r_i}
			      \right|_{\sigma(t)} \equiv (\dot{\sigma}_1, \dots,
			      \dot{\sigma}_d)$.
	\end{itemize}
\end{frame}

%====================================================================
\section{An Application}

\begin{frame}{A map with vanishing differential is constant}
	\begin{theorem}
		Let $\Psi$ be a $C^\infty$ mapping of the connected manifold $M$
		into $N$. If $d\Psi|_p \equiv 0$ for all $p \in M$, then $\Psi$ is
		a constant map.
	\end{theorem}
	\begin{proof}
		Let $q \in \Psi(M)$. The set $\Psi^{-1}(\{q\})$ is closed; we show
		it is open. Let $m \in \Psi^{-1}(\{q\})$, and choose coordinate
		systems $(U, x_1, \dots, x_d)$, $(V, y_1, \dots, y_e)$ about $m$
		and $q$ with $\Psi(U) \subset V$. Then for all $p \in U$,
		\[
			0 = d\Psi\left( \frac{\partial}{\partial x_j} \right)
			= \sum_{i=1}^{e} \left. \frac{\partial (y_i \circ \Psi)}
			{\partial x_j} \right|_{p}
			\left. \frac{\partial}{\partial y_i} \right|_{\Psi(p)}
			\ \Longrightarrow\
			\left. \frac{\partial (y_i \circ \Psi)}{\partial x_j}
			\right|_{p} = 0 \quad \forall\, i, j.
		\]
		So the $y_i \circ \Psi$ are constant on $U$, i.e.\ $\Psi(U) = \{q\}$,
		and $\Psi^{-1}(\{q\})$ is open. \qedhere
	\end{proof}
\end{frame}

%====================================================================
\section{Tangent and Cotangent Bundles}

\begin{frame}{Tangent and cotangent bundles as sets}
	Let $M$ be a $C^\infty$ manifold with differentiable structure
	$\mathcal{F}$.
	\begin{definition}[Tangent \& cotangent bundle]
		\[
			TM = \bigcup_{p \in M} T_p M,
			\qquad
			T^* M = \bigcup_{p \in M} T_p^* M
		\]
		as two sets of tangent / cotangent vectors.
	\end{definition}
	\begin{definition}[Natural projection]
		\[
			\pi : TM \longrightarrow M, \qquad
			\pi^* : T^* M \longrightarrow M,
		\]
		sending $v$ (at $p$) $\mapsto p$ and $\omega$ (at $p$) $\mapsto p$,
		where $v \in T_p M$, $\omega \in T_p^* M$.
	\end{definition}
\end{frame}

\begin{frame}{Local trivializations}
	\begin{definition}[Local trivialization]
		Let $(U, \varphi)$ be a local coordinate system in $\mathcal{F}$,
		with coordinate functions $x_1, \dots, x_d$. Define maps
		\[
			\widetilde{\varphi} : \pi^{-1}(U) \longrightarrow \R^{2d},
			\qquad
			\widetilde{\varphi}^* : (\pi^*)^{-1}(U) \longrightarrow \R^{2d}
		\]
		by
		\[
			\widetilde{\varphi}(w)
			= \bigl( \varphi(p),\ a_1, \dots, a_d \bigr),
		\]
		if $w \in T_p M$ can be written as
		$w = \sum_{i=1}^{d} a_i \left. \frac{\partial}{\partial x_i} \right|_{p}
			\in T_p M \subset \pi^{-1}(U)$; and similarly
		$\widetilde{\varphi}^*(\omega)
			= \bigl( \varphi(p),\ b_1, \dots, b_d \bigr)$
		if $\omega = \sum_{j=1}^{d} b_j\, dx_j \in T_p^* M$.
	\end{definition}
	\begin{block}{Note}
		Both $\widetilde{\varphi}$ and $\widetilde{\varphi}^*$ are
		one-to-one maps onto open subsets
		$\varphi(U) \times \R^d \subset \R^{2d}$.
	\end{block}
\end{frame}

\begin{frame}{The topology on $TM$}
	\emph{Outline of the construction of the $C^\infty$ manifold
		structure on $TM$.}
	\begin{enumerate}
		\item $TM$ is a topological space, which is Hausdorff, second
		      countable and locally compact.
	\end{enumerate}
	\begin{itemize}
		\item The collection
		      \[
			      \mathcal{B} \triangleq \bigl\{ \widetilde{\varphi}^{-1}(W)
			      \bigm| W \subseteq \R^{2d} \ \text{open},\ W \neq
			      \varnothing,\ \forall\, (U, \varphi) \in \mathcal{F} \bigr\}
		      \]
		      forms a basis for a topology on $TM$.
		\item The topology $\mathcal{V}$ generated by $\mathcal{B}$ makes
		      $(TM, \mathcal{V})$ Hausdorff, second countable and locally
		      compact.
		\item $TM$ is a local Euclidean space because of the local
		      trivialization $\widetilde{\varphi}$ (which provides the
		      homeomorphism).
	\end{itemize}
\end{frame}

\begin{frame}{The differentiable structure on $TM$}
	The compatible differentiable structure $\widetilde{\mathcal{F}}$ of
	$TM$ is the maximal collection of local coordinate systems containing
	\[
		\widetilde{\mathcal{F}} \supset
		\bigl\{ \bigl( \pi^{-1}(U), \widetilde{\varphi} \bigr)
		\bigm| (U, \varphi) \in \mathcal{F} \bigr\}.
	\]
	\begin{exercise}
		It is necessary to show that, given two local coordinate systems
		$\bigl( \pi^{-1}(U), \widetilde{\varphi} \bigr)$ and
		$\bigl( \pi^{-1}(V), \widetilde{\psi} \bigr)$, the transition map
		\[
			\widetilde{\psi} \circ \widetilde{\varphi}^{-1}
		\]
		is a $C^\infty$ map. (Check it!)
	\end{exercise}
	\bigskip
	$TM$ and $T^*M$ with these differentiable structures are called the
	\emph{tangent bundle} and \emph{cotangent bundle}, respectively. It is
	sometimes convenient to write points of $TM$ as pairs $(p, v)$, where
	$p \in M$ and $v \in T_p M$.
\end{frame}

\begin{frame}{Differentials of bundle maps}
	If $\Psi : M \to N$ is a $C^\infty$ map, then the differential of
	$\Psi$ defines a mapping of the tangent bundles
	\[
		d\Psi = \Psi_* : TM \longrightarrow TN,
	\]
	where
	\[
		d\Psi(p, w) = \bigl( \Psi(p),\ d\Psi(w) \bigr).
	\]
	\begin{exercise}
		Check that this is a $C^\infty$ map.
	\end{exercise}
\end{frame}

\end{document}
